\section{Hardware realisation}

The hardware platform is an ESP32-WROOM module connected to an L298N H-bridge.
The L298N introduces an approximate voltage drop of \(\SI{4}{V}\), so with a
\(\SI{12}{V}\) supply the motor sees about \(\SI{8}{V}\). The motor is equipped
with a quadrature encoder. The sampled controller uses the period
\(T=\SI{0.08}{s}\), \(P(0)=100I\),
\(x_e(0)=\begin{bmatrix}-0.5&1.0\end{bmatrix}^{T}\), forgetting factor
\(\alpha=0.98\), desired polynomial \(q(z)=z^2-0.5z+0.06\), and a command
saturated at \(\pm255\).

The controller is implemented by automatic code generation from Simulink using
the Embedded Real-Time (ERT) target, the ESP32 target support, and External
Mode for Monitor and Tune operation. Table~\ref{tab:esp32-pin-assignment}
gives the pin assignment.

\begin{scripttable}
  \centering
  \caption{ESP32 pin assignment for the DC motor test rig.}
  \label{tab:esp32-pin-assignment}
  \begin{tabular}{ll}
    \toprule
    Signal & ESP32 pin \\
    \midrule
    PWM & GPIO 14 \\
    IN1 & GPIO 26 \\
    IN2 & GPIO 27 \\
    Encoder A & GPIO 32 \\
    Encoder B & GPIO 33 \\
    \bottomrule
  \end{tabular}
\end{scripttable}

The encoder resolution was determined by a one-revolution test as
\(44\) counts per revolution, corresponding to \(11\) encoder impulses with
x4 quadrature edge decoding. The speed conversion
\[
  \omega = \frac{2\pi \Delta N}{T\,\mathrm{CPR}}
\]
is therefore calibrated for \(\mathrm{CPR}=44\). The measured full-duty speed
of approximately \(\SI{780}{rad/s}\) corresponds to about \(\SI{7450}{rpm}\),
consistent with the gearbox-free 25GA371 speed range of roughly
\(\SIrange{3000}{8000}{rpm}\) and independently supporting the encoder scaling.
